\section{Background}\label{sec:background}

\subsection{Multiresolution Analysis in MADNESS}\label{sec:mra}

MADNESS~\cite{harrison2016madness} represents three-dimensional functions in a discontinuous multiwavelet basis built on an adaptive octree.
The simulation cell $[0,1)^3$ is recursively bisected along each axis to produce cubic boxes at integer levels $n = 0, 1, 2, \ldots$, with each box at level $n$ identified by a translation index $\mathbf{l} \in \{0, \ldots, 2^n - 1\}^3$.
Within each box, the function is expanded in a tensor product of Legendre scaling polynomials of order $k$ (we use $k = 8$ throughout), giving $k^3 = 512$ scaling coefficients (s-coefficients) per node.

The connection between consecutive levels is mediated by Alpert's two-scale relation~\cite{alpert1993twoscale}.
A parent node's s-coefficients can be decomposed into the s-coefficients of its eight children plus a set of wavelet difference coefficients (d-coefficients) that encode the detail lost in the coarsening.
Refinement is controlled by the d-coefficient norm: if $\|\mathbf{d}\| > \epsilon$ for a prescribed threshold $\epsilon$, the node is subdivided into eight children and the representation is extended to the next level.
Nodes whose d-coefficients fall below $\epsilon$ become leaves.
This adaptive refinement concentrates computational effort near nuclei and in bonding regions while keeping the representation coarse in smooth, slowly varying parts of space.

\subsection{SCF Iteration Mechanics}\label{sec:scf}

Kohn-Sham DFT~\cite{kohn1965self} reduces the many-electron problem to a set of single-particle eigenvalue equations in which the effective potential depends on the electron density $\rho(\mathbf{r})$.
Because the potential is a functional of $\rho$, the equations must be solved self-consistently: an initial density defines a potential, the eigenvalue equations are solved for new orbitals, a new density is constructed, and the cycle repeats until $\rho$ changes by less than a convergence threshold.

MADNESS uses the promolecular density $\rho_0(\mathbf{r}) = \sum_A \rho_A^\mathrm{atom}(\mathbf{r})$ as its initial guess, where $\rho_A^\mathrm{atom}$ is the spherically symmetric ground-state density of isolated atom $A$.
This superposition is inexpensive to construct and places charge in roughly the right spatial regions, but it contains no information about chemical bonding, charge redistribution, or exchange-correlation effects.
The distance between $\rho_0$ and the converged $\rho$ determines how many SCF iterations are needed.

MADNESS employs a multi-protocol convergence strategy~\cite{harrison2016madness}.
The calculation begins at a coarse truncation threshold (typically $\epsilon = 10^{-4}$), producing a shallow tree with fewer nodes.
Once the density converges at this threshold, the tree is refined to a tighter threshold ($\epsilon = 10^{-6}$ or beyond) and iteration continues.
The first (coarsest) protocol dominates the total iteration count because it starts from $\rho_0$, the density furthest from the converged solution.
Subsequent protocols start from an already-converged density at the previous threshold and typically converge in one or two iterations.
Any ML model that aims to reduce SCF cost must therefore improve the density at coarse levels, where the first protocol operates and where the gap between $\rho_0$ and $\rho$ is largest.
